\chapter{Equilibrium}
\begin{definition}{Moment}{moment}
	The moment of a force about a pivot is the product of (1) the magnitude of the force, and
	(2) the distance of the \it{line of action} of the force from the pivot.
	\[\tau = Fd_\perp\]
\end{definition}
\begin{definition}{Couple}{couple}
	Two \it{parallel} forces of equal magnitude but opposite direction, whose \it{lines of action}
	do not coincide. Where \(d\) is the perpendicular distance between the forces,
	\[\tau_\text{couple} = Fd\]
\end{definition}
\begin{definition}{Equilibrium}{equilibrium}
	For a body to be in equilibrium, it must satisfy two conditions.
	\begin{itemize}
		\item The \it{net external force} acting on the body must be zero---translational equilibrium.
		      \[\sum \va{F} = 0 \implies \sum F_x = 0, \sum F_y = 0\]
		\item The \it{net torque} acting on the body \it{about any point} must be zero---rotational equilibrium.
		      \[\sum \tau = 0 \implies \sum\tau_\text{ACW}=\sum\tau_\text{CW}\quad\text{for any pt.}\]
	\end{itemize}
\end{definition}
\begin{definition}{Centre of gravity}{cg}
	The point at which the weight of the body appears to act.
\end{definition}
For a body whose mass is uniformly distributed, its c.g.~coincides with its \it{geometric centre}.